\clearpage
\VerdierPhysicalPageAnchor{0245}%
\begingroup\small\setlength{\parskip}{0pt}\setlength{\emergencystretch}{1.5em}
\setlength{\abovedisplayskip}{4pt plus 2pt minus 2pt}%
\setlength{\belowdisplayskip}{4pt plus 2pt minus 2pt}%

\VerdierSemanticAnchor{ch4d-4-6-3-b}
\noindent\textit{b)
$Q\colon K(\mathcal A)\longrightarrow D(\mathcal A)$ を商関手とする。複体
$I^\bullet$ は $Q$-右自由である（第 II 章、2.3.3）。すなわち任意の複体
$Y$ に対し、標準射}
\[
\operatorname{Hom}_{K(\mathcal A)}(Y,I^\bullet)
\longrightarrow\operatorname{Hom}_{D(\mathcal A)}(Y,I^\bullet)
\]
\textit{は同型である。}

\VerdierSemanticAnchor{ch4d-4-6-3-c}
\noindent\textit{c)}
\[
H^*(I^\bullet)\colon\cdots\longrightarrow H^p(I^\bullet)
\longrightarrow H^{p+1}(I^\bullet)\longrightarrow\cdots
\]
\textit{を、成分が $H^p(I^\bullet)$、$p\in\mathbb Z$ で、微分が零である
複体とする。コホモロジー対象上で恒等射を誘導する $K(\mathcal A)$ の同型
$I^\bullet\longrightarrow H^*(I^\bullet)$ が一意に存在する。}

\VerdierSemanticAnchor{ch4d-4-6-3-d}
\noindent\textit{d) 標準射
$H^*(I^\bullet)\longrightarrow H^p(I^\bullet)[-p]$ は、$K(\mathcal A)$ と
$D(\mathcal A)$ において同型を定める：}
\[
H^*(I^\bullet)\xrightarrow{\sim}
\prod_{p\in\mathbb Z}H^p(I^\bullet)[-p].
\]

\VerdierSemanticAnchor{ch4d-4-6-3-proof-abcd}
a) が b) を含意することを示す。実際、任意の非輪状複体 $Y$ に対して
$\operatorname{Hom}_{K(\mathcal A)}(Y,I^\bullet)=0$ であることを示せば十分で
あり（第 II 章、2.3.3）、これは a) から従う。a) が c) を含意することを
示す。実際、a) の同型において複体 $Y$ を $I^\bullet$ で、複体
$I^\bullet$ を $H^*(I^\bullet)$ で置き換えればよい。同様に、a) の同型で
複体 $I^\bullet$ を複体 $H^*(I^\bullet)$ で置き換えると、a) が d) を
含意することが分かる。a) を示す。定義から直ちに、複体 $I^\bullet$ は
次の形の複体と同型である：

\VerdierSemanticAnchor{ch4d-4-6-3-proof-a-decomposition}
\[
\resizebox{0.98\linewidth}{!}{$
\displaystyle
\cdots\longrightarrow
J^{n-2}\oplus H^{n-1}\oplus J^{n-1}
\xrightarrow{\left(\begin{smallmatrix}
0&0&\mathrm{id}\\0&0&0\\0&0&0
\end{smallmatrix}\right)}
J^{n-1}\oplus H^n\oplus J^n
\xrightarrow{\left(\begin{smallmatrix}
0&0&\mathrm{id}\\0&0&0\\0&0&0
\end{smallmatrix}\right)}
J^n\oplus H^{n+1}\oplus J^{n+1}
\longrightarrow\cdots
$}.
\]

\VerdierSemanticAnchor{ch4d-4-6-3-proof-a-homotopy}
したがって、$I^\bullet$ と複体 $H^*(I^\bullet)$ の間にはホモトピー同値が
存在する。a) を示すには $I^\bullet=H^*(I^\bullet)$ と仮定してよい。
このとき明らかに同型

\VerdierSemanticAnchor{ch4d-4-6-3-proof-a-hom}
\[
\operatorname{Hom}_{K(\mathcal A)}\bigl(Y,H^*(I^\bullet)\bigr)
\xrightarrow{\sim}
\prod_{p\in\mathbb Z}
\operatorname{Hom}_{K(\mathcal A)}
\bigl(Y,H^p(I^\bullet)[-p]\bigr).
\]
がある。さらに、対象 $H^p(I^\bullet)$ は入射対象なので、標準射
\[
\operatorname{Hom}_{K(\mathcal A)}
\bigl(Y,H^p(I^\bullet)[-p]\bigr)
\longrightarrow
\operatorname{Hom}_{\mathcal A}\bigl(H^p(Y),H^p(I^\bullet)\bigr)
\]
は同型である。
\endgroup
